\documentclass[10pt]{article}
\usepackage{amsmath,amssymb,mathtools}
\usepackage[margin=0.65in]{geometry}
\allowdisplaybreaks
\begin{document}

\section*{Waste collection inventory-routing problem}
\begin{equation}
\begin{aligned}
(\mathrm{WCIRP})\qquad &\min\quad z=\mathrm{RC}+\mathrm{ESOEDC}+\mathrm{ERFC}\\
&\mathrm{s.t.}\quad\text{constraints (18)--(47), expanded below, and (55).}
\end{aligned}\tag{54}
\end{equation}
The selected formulation is Section 6.3 of Markov, Bierlaire, Cordeau, Maknoon and Varone, \emph{A Unified Framework for Rich Routing Problems with Stochastic Demands}, February 27, 2018, physical PDF pages 18--19. Collection is interpreted as distribution of empty space: $\mathcal P$ is the set of waste containers and $\mathcal D$ the set of disposal facilities. No alternative chance constraints are imposed.

\subsection*{Subject to: inherited constraints}
\begin{align}
\sum_{o'\in\mathcal O'_{kt}}\sum_{j\in\mathcal N_{kt}}x_{o'jkt}
&=\sum_{i\in\mathcal N_{kt}}\sum_{o''\in\mathcal O''_{kt}}x_{io''kt},
&&\forall t\in\mathcal T,\ k\in\mathcal K.\tag{18}\\
\sum_{i\in\mathcal N_{kt}}x_{io'kt}&=0,
&&\forall t\in\mathcal T,\ k\in\mathcal K,\ o'\in\mathcal O'_{kt}.\tag{19}\\
\sum_{j\in\mathcal N_{kt}}x_{o''jkt}&=0,
&&\forall t\in\mathcal T,\ k\in\mathcal K,\ o''\in\mathcal O''_{kt}.\tag{20}\\
\sum_{i\in\mathcal N_{kt}}x_{iokt}
&=\sum_{j\in\mathcal N_{k(t+1)}}x_{ojk(t+1)},
&&\forall t\in\mathcal T,\ k\in\mathcal K,\ o\in\mathcal O''_{kt}\cap\mathcal O'_{k(t+1)}.\tag{21}\\
 y_{ikt}&=\sum_{j\in\mathcal N_{kt}}x_{ijkt},
&&\forall t\in\mathcal T,\ k\in\mathcal K,\ i\in\mathcal N_{kt}\setminus\mathcal O''_{kt}.\tag{22}\\
 y_{jkt}&=\sum_{i\in\mathcal N_{kt}}x_{ijkt},
&&\forall t\in\mathcal T,\ k\in\mathcal K,\ j\in\mathcal O''_{kt}.\tag{23}\\
\sum_{k\in\mathcal K}y_{ikt}&\leqslant1,
&&\forall t\in\mathcal T,\ i\in\mathcal P.\tag{24}\\
 y_{ikt}&\leqslant\alpha_{ikt},
&&\forall t\in\mathcal T,\ k\in\mathcal K,\ i\in\mathcal P\cup\mathcal D.\tag{25}\\
\sum_{i\in\mathcal N_{kt}}x_{ijkt}&=\sum_{i\in\mathcal N_{kt}}x_{jikt},
&&\forall t\in\mathcal T,\ k\in\mathcal K,\ j\in\mathcal P\cup\mathcal D.\tag{26}\\
\sum_{r\in\mathcal C_i}c_{ir}&=1,
&&\forall i\in\mathcal P.\tag{27}\\
\sum_{k\in\mathcal K}y_{ikt}-\sum_{r\in\mathcal C_i}\alpha_{rt}c_{ir}&=0,
&&\forall t\in\mathcal T,\ i\in\mathcal P.\tag{28}\\
 I_{it}&=I_{i(t-1)}+\sum_{k\in\mathcal K}q_{ik(t-1)}-\mathbb E(\rho_{i(t-1)}),
&&\forall t\in\mathcal T^+,\ i\in\mathcal P.\tag{29}\\
 I_{it}&\geqslant0,
&&\forall t\in\mathcal T^+,\ i\in\mathcal P.\tag{30}\\
 -I_{i0}&\leqslant\omega_i\sum_{k\in\mathcal K}y_{ik0},
&&\forall i\in\mathcal P.\tag{31}\\
\sum_{k\in\mathcal K}y_{ikt}-\sum_{r\in\mathcal L_i}\ell_{irt}&=0,
&&\forall t\in\mathcal T,\ i\in\mathcal P.\tag{32}\\
 q_{ikt}&\geqslant\sum_{r\in\mathcal L_i}r\ell_{irt}-I_{it}-M(1-y_{ikt}),
&&\forall t\in\mathcal T,\ k\in\mathcal K,\ i\in\mathcal P.\tag{33}\\
 q_{ikt}&\leqslant\sum_{r\in\mathcal L_i}r\ell_{irt}-I_{it}+\omega_i(1-y_{ikt}),
&&\forall t\in\mathcal T,\ k\in\mathcal K,\ i\in\mathcal P.\tag{34}\\
 q_{ikt}&\leqslant My_{ikt},
&&\forall t\in\mathcal T,\ k\in\mathcal K,\ i\in\mathcal P.\tag{35}\\
 q_{ikt}&\leqslant Q_{ikt}\leqslant\Omega_k,
&&\forall t\in\mathcal T,\ k\in\mathcal K,\ i\in\mathcal P.\tag{36}\\
 Q_{ikt}&=0,
&&\forall t\in\mathcal T,\ k\in\mathcal K,\ i\in\mathcal D.\tag{37}\\
 Q_{ikt}+q_{jkt}&\leqslant Q_{jkt}+\Omega_k(1-x_{ijkt}),
&&\forall t\in\mathcal T,\ k\in\mathcal K,\ i\in\mathcal N_{kt},\ j\in\mathcal N_{kt}\setminus\mathcal D.\tag{38}\\
 Q_{o'k(t+1)}&\geqslant Q_{o''kt},
&&\forall t\in\mathcal T,\ k\in\mathcal K,\ o'\in\mathcal O'_{kt},\ o''\in\mathcal O''_{kt}.\tag{39}\\
 S_{ikt}+\delta_i+\tau_{ijkt}
&\leqslant S_{jkt}+(\mu_i+\delta_i+\tau_{ijkt})(1-x_{ijkt}),
&&\forall t\in\mathcal T,\ k\in\mathcal K,\ i\in\mathcal N_{kt},\ j\in\mathcal N_{kt}.\tag{40}\\
\lambda_i y_{ikt}&\leqslant S_{ikt}\leqslant\mu_i y_{ikt},
&&\forall t\in\mathcal T,\ k\in\mathcal K,\ i\in\mathcal N_{kt}.\tag{41}\\
\underline b_{kt}&\leqslant\sum_{o''\in\mathcal O''_{kt}}S_{o''kt}
-\sum_{o'\in\mathcal O'_{kt}}S_{o'kt}\leqslant\bar b_{kt},
&&\forall t\in\mathcal T,\ k\in\mathcal K.\tag{42}\\
\bar b_{kt}&\leqslant\mathrm H_{kt}z_{kt},
&&\forall t\in\mathcal T,\ k\in\mathcal K.\tag{43}\\
\underline B&\leqslant\sum_{t\in\mathcal T}\underline b_{kt},
&&\forall k\in\mathcal K.\tag{44}\\
\bar B&\geqslant\sum_{t\in\mathcal T}\bar b_{kt},
&&\forall k\in\mathcal K.\tag{45}
\end{align}
\begin{align}
 x_{ijkt},y_{ikt},z_{kt},c_{ir'},\ell_{ir''t}&\in\{0,1\},
&&\forall t\in\mathcal T,\ k\in\mathcal K,\ i,j\in\mathcal N_{kt},\ r'\in\mathcal C_i,\ r''\in\mathcal L_i.\tag{46}\\
 q_{ikt},Q_{ikt},I_{it},S_{ikt},\underline b_{kt},\bar b_{kt},\underline B,\bar B&\geqslant0,
&&\forall t\in\mathcal T,\ k\in\mathcal K,\ i\in\mathcal N_{kt}.\tag{47}
\end{align}
The big-$M$ in (33) and (35) has the point- and period-dependent values specified in the prose on PDF page 16:
\[
 M=\begin{cases}2\omega_i,&t=0,\\\omega_i,&\text{otherwise}.\end{cases}
\]
\subsection*{Additional WCIRP constraint}
\begin{equation}
 Q_{o''kt}=0,\qquad\forall t\in\mathcal T,\ k\in\mathcal K,\ o''\in\mathcal O''_{kt}.\tag{55}
\end{equation}
The authors state that (55), together with (37)--(38), forces a disposal-facility visit immediately before the destination depot. This statement does not add an arc restriction to the printed model.

\section*{Sets, fixed data, and decision roles}
The definitions are from Section 3 and Table 1 (PDF pages 6--8) and Section 5 (page 13).
\begin{gather*}
\mathcal T=\{0,\ldots,u\},\qquad
\mathcal T^+=\{1,\ldots,u,u+1\},\\
\mathcal O'_{kt}\subseteq\mathcal O,\qquad
\mathcal O''_{kt}\subseteq\mathcal O,\\
\mathcal N_{kt}=\mathcal O'_{kt}\cup\mathcal O''_{kt}\cup\mathcal P\cup\mathcal D,\\
\mathcal A_{kt}=\{(i,j):\ \forall i,j\in\mathcal N_{kt},\ i\ne j\},\qquad
G_{kt}(\mathcal N_{kt},\mathcal A_{kt}),\quad k\in\mathcal K,\ t\in\mathcal T.
\end{gather*}
$\mathcal K$ is a fixed heterogeneous fleet; $\mathcal O$ contains all origin and destination depots, and $\mathcal O'_{kt}$ and $\mathcal O''_{kt}$ are the origin and destination choices for vehicle $k$ in period $t$. The set $\mathcal D$ contains sufficiently many replications of each supply point to allow repeated visits by a vehicle within a period. The distance matrix is asymmetric. $\pi_{ij}$ is the length of arc $(i,j)\in\mathcal A_{kt}$; $\tau_{ijkt}$ is its vehicle- and period-specific travel time.

For $k\in\mathcal K$, $\varphi_k$ is the per-period deployment cost, $\beta_k$ the unit-distance running cost, $\theta_k$ the unit-time running cost, and $\Omega_k$ the capacity. For $i\in\mathcal O\cup\mathcal P\cup\mathcal D$, $[\lambda_i,\mu_i]$ is the start-of-service time window and $\delta_i$ the service duration. The following fixed-data conditions are stated:
\[
\delta_i=0\quad(i\in\mathcal O),\qquad
\mathcal O''_{kt}\cap\mathcal O'_{k(t+1)}\ne\varnothing.
\]
A vehicle arriving before $\lambda_i$ waits; service may not start after $\mu_i$. $\mathrm H_{kt}$ is its maximum tour duration. If $\mathrm H_{kt}=0$, the vehicle is unavailable, and the depot sets must permit idling:
\[
\mathrm H_{kt}=0\ \Longrightarrow\
\exists o'\in\mathcal O'_{kt},\ o''\in\mathcal O''_{kt}\ \text{such that }\pi_{o'o''}=0.
\]
The depot compatibility and idling conditions are stated for the vehicle-period depot sets in Section 3; no terminal-period convention is added here.

For $i\in\mathcal P$, $\omega_i$ is inventory capacity. $\mathcal C_i$ contains allowable visit-period combinations, possibly with different frequencies. The fixed binary $\alpha_{rt}$ is one if period $t$ belongs to combination $r\in\mathcal C_i$, and zero otherwise. The separate fixed binary $\alpha_{ikt}$ is one if point $i\in\mathcal P\cup\mathcal D$ is accessible to vehicle $k$ in period $t$, and zero otherwise; it also expresses continuity of service. Demand points are served at most once per period. The general notation introduces a minimum visit count $\nu_i$, but the selected WCIRP block contains (27)--(28), not an additional minimum-count inequality.

The selected inventory policy is the discretized maximum-level policy. $\mathcal L_i$ is the given set of allowable discrete inventory levels after regular delivery. The source does not prescribe a numerical discretization. The order-up-to policy is the special case $\mathcal L_i=\{\omega_i\}$ for every $i\in\mathcal P$; this special case is not imposed here. Emergency deliveries always use the order-up-to policy.

The fixed monetary parameters $\chi_i$ and $\zeta_i$ are, respectively, the stock-out and emergency-delivery costs; for collection they price overflow and emergency collection. The route-failure cost multiplier satisfies
\[
\psi\in[0,1].
\]
$C_{\mathcal S}$ is the average routing cost of going from the demand points of trip $\mathcal S$ to their nearest supply point and back. It remains the source-defined abstract cost; the paper gives no further algebraic formula for it.

The binary decision roles are: $x_{ijkt}=1$ for traversal of arc $(i,j)$ by vehicle $k$ in period $t$; $y_{ikt}=1$ for a visit to point $i$; $z_{kt}=1$ for use of vehicle $k$; $c_{ir}=1$ for assignment of visit combination $r$ to point $i$; and $\ell_{irt}=1$ for selection of discrete level $r\in\mathcal L_i$ at point $i$ in period $t$. Each is zero otherwise.

The continuous variables are $q_{ikt}$, the expected delivery quantity to demand point $i$; $Q_{ikt}$, the expected cumulative quantity delivered by vehicle $k$ arriving at point $i$; $I_{it}$, the expected inventory at the start of period $t$, before delivery; and $S_{ikt}$, the start-of-service time. The variables $\underline b_{kt},\bar b_{kt}$ bound tour duration, while $\underline B,\bar B$ bound each vehicle's total workload over the horizon. The selected objective does not price workload balancing, but (42)--(47) remain inherited. The known initial inventory $I_{i0}$ is fixed data, not a freely chosen initial state.

\section*{Cost definitions used by (54)}
\begin{equation}
\mathrm{RC}=\sum_{t\in\mathcal T}\sum_{k\in\mathcal K}
\left(\varphi_kz_{kt}+\beta_k\sum_{i\in\mathcal N_{kt}}\sum_{j\in\mathcal N_{kt}}\pi_{ij}x_{ijkt}
+\theta_k\left(\sum_{o''\in\mathcal O''_{kt}}S_{o''kt}-\sum_{o'\in\mathcal O'_{kt}}S_{o'kt}\right)\right).
\tag{13}
\end{equation}
\begin{equation}
\mathrm{ESOEDC}=\sum_{t\in\mathcal T\cup\mathcal T^+}\sum_{i\in\mathcal P}
\left(\chi_i+\zeta_i-\zeta_i\sum_{k\in\mathcal K}y_{ikt}\right)\mathsf p^{\mathrm{DP}}_{it}.
\tag{15}
\end{equation}
\begin{equation}
\mathrm{ERFC}=\sum_{k\in\mathcal K}\sum_{\mathcal S\in\mathfrak S_k}
\psi C_{\mathcal S}\mathsf p^{\mathrm{RF}}_{\mathcal S,k}.
\tag{16}
\end{equation}
Here $\mathfrak S_k$ is the solution-dependent set of supply-point-delimited trips of vehicle $k$, $\mathcal S$ is a particular trip, and $\mathcal S_t\in\mathcal S$ is the set of its demand points visited in period $t$. Supply points delimiting a trip can be visited in different periods. These sets are distinct from the start-of-service variable $S_{ikt}$.

In a stock-out period, (15) charges $\chi_i$ when a regular visit occurs and $\chi_i+\zeta_i$ otherwise. Initial stock-outs are charged even though the initial inventory is known. The scope $\mathcal T\cup\mathcal T^+$ includes the first period after the horizon because its stock-out probability is determined by horizon decisions.

\paragraph{Cost-model qualification (Assumption 4, PDF pages 14--15).}
All objective terms except ESOEDC ignore the cost effect of demand points stocking out earlier than expected. A reaction policy means changes later in the horizon in response to recourse, ranging from doing nothing to re-optimization; it is not an additional decision block here. In the absence of inventory holding cost, Proposition 2 states that this objective overestimates real cost. In particular, emergency delivery can reduce a later trip's delivery quantity and route-failure probability, an effect omitted by (8) and (16).

\section*{Demand, inventory action sequence, and stock-out probabilities}
Demands are the only stochastic inputs; all other parameters are deterministic. They may be non-stationary. For $i\in\mathcal P$ and $t\in\mathcal T$,
\begin{equation}
\rho_{it}=\mathbb E(\rho_{it})+\varepsilon_{it}.
\tag{1}
\end{equation}
The source writes the vector of errors as
\begin{equation}
\varepsilon=(\varepsilon_{11},\ldots,\varepsilon_{1|\mathcal T|},\varepsilon_{21},\ldots,\varepsilon_{|\mathcal P||\mathcal T|}).
\tag{2}
\end{equation}
\[
\varepsilon\sim\Phi,\qquad\operatorname{var}(\varepsilon)=K.
\]
Here $K$ is a covariance structure, not the fleet $\mathcal K$. Definition 1 requires the forecasting model to provide $\mathbb E(\rho_{it})$ for every $t\in\mathcal T,i\in\mathcal P$ and the joint error distribution $\Phi$. The latter may be theoretical or empirical but must be simulatable. No normality restriction is imposed in this formulation.

\paragraph{Initial state and timing (Definitions 2--3).}
For every $i\in\mathcal P$, $I_{i0}$ is observed and known with certainty and may be positive, zero, or negative. Consequently its initial stock-out probability is either zero or one. In each period $t\in\mathcal T$, deliveries occur before demand realizations. A regular delivery in period $t$ means $y_{ikt}=1$ for some $k\in\mathcal K$. The random state flag $\sigma_{it}$ is one for stock-out and zero otherwise; it is not a routing decision. An emergency delivery occurs when
\[
\sigma_{it}=1\quad\text{and}\quad y_{ikt}=0\ \text{for all }k\in\mathcal K,
\]
and restores inventory to $\omega_i$. The limited-backorder policy requires delivery in the same period as a stock-out.

Let $\Lambda_{ig}$ denote inventory after a regular delivery to point $i$ in period $g$. Its discrete choice is represented by $\ell_{irg}$:
\[
\Lambda_{ig}\in\mathcal L_i.
\]
When no delivery to point $i$ occurs during the planning horizon, the source specifies $\Lambda_{i0}=I_{i0}$. In (7), period zero is also returned when no earlier regular delivery exists.

\paragraph{Assumption 1.}
A regular delivery to $i\in\mathcal P$ in $t\in\mathcal T$ discards any inventory above its chosen level $\Lambda_{it}$. This is the selected rule for handling a realized inventory above a discretized maximum-level target; neither skipping delivery nor picking up excess inventory replaces it.

For a regular delivery to point $i$ in period $g\in\mathcal T$, the four stock-out transition probabilities are:
\begin{equation}
\mathbb P(\Lambda_{ig}-\rho_{ig}\leqslant0)
=\mathbb P\bigl(\varepsilon_{ig}\geqslant\Lambda_{ig}-\mathbb E(\rho_{ig})\bigr).
\tag{3}
\end{equation}
\begin{equation}
\begin{aligned}
&\mathbb P\left(\Lambda_{ig}-\sum_{t=g}^{h}\rho_{it}\leqslant0\ \middle|\ \Lambda_{ig}-\sum_{t=g}^{h-1}\rho_{it}>0\right)\\
&\quad=\mathbb P\left(\sum_{t=g}^{h}\varepsilon_{it}\geqslant\Lambda_{ig}-\sum_{t=g}^{h}\mathbb E(\rho_{it})\ \middle|\ \sum_{t=g}^{h-1}\varepsilon_{it}<\Lambda_{ig}-\sum_{t=g}^{h-1}\mathbb E(\rho_{it})\right),\qquad\forall h>g.
\end{aligned}\tag{4}
\end{equation}
\begin{equation}
\mathbb P(\omega_i-\rho_{ig'}\leqslant0)
=\mathbb P\bigl(\varepsilon_{ig'}\geqslant\omega_i-\mathbb E(\rho_{ig'})\bigr),\qquad\forall g'>g.
\tag{5}
\end{equation}
\begin{equation}
\begin{aligned}
&\mathbb P\left(\omega_i-\sum_{t=g'}^{h}\rho_{it}\leqslant0\ \middle|\ \omega_i-\sum_{t=g'}^{h-1}\rho_{it}>0\right)\\
&\quad=\mathbb P\left(\sum_{t=g'}^{h}\varepsilon_{it}\geqslant\omega_i-\sum_{t=g'}^{h}\mathbb E(\rho_{it})\ \middle|\ \sum_{t=g'}^{h-1}\varepsilon_{it}<\omega_i-\sum_{t=g'}^{h-1}\mathbb E(\rho_{it})\right),\qquad\forall h>g'>g.
\end{aligned}\tag{6}
\end{equation}
Equations (3)--(4) concern transitions from a regular delivery; (5)--(6) concern transitions from a subsequent stock-out followed by an emergency delivery to capacity. The binary tree in Figure 1 (PDF page 10) assigns these probabilities to branches entering state one and the corresponding positive-inventory probabilities to branches entering state zero. For example, the printed state-zero branches include
\begin{gather*}
\mathbb P(\Lambda_{ig}-\rho_{ig}>0),\\
\mathbb P(\Lambda_{ig}-\rho_{ig}-\rho_{i(g+1)}>0\mid\Lambda_{ig}-\rho_{ig}>0),\\
\mathbb P(\omega_i-\rho_{i(g+1)}>0),\\
\mathbb P(\omega_i-\rho_{i(g+1)}-\rho_{i(g+2)}>0\mid\omega_i-\rho_{i(g+1)}>0).
\end{gather*}
The stock-out probability at $t>g$ is the sum of the probabilities of all tree paths from the root inventory $\Lambda_{ig}$ in period $g$ to state $\sigma_{it}=1$. The probability in period $g$ comes from the preceding tree, or is zero or one when $g=0$. The resulting definition is
\begin{equation}
\mathsf p^{\mathrm{DP}}_{it}
=\mathbb P\left(\sigma_{it}=1\ \middle|\ \Lambda_{im}:\ m=\max(0,g\in\mathcal T:g<t:\exists k\in\mathcal K:y_{ikg}=1)\right).
\tag{7}
\end{equation}
The index $m$ identifies the most recent earlier regular delivery, with zero as the printed fallback. Under the discretized maximum-level policy these probabilities can be precomputed for each $r\in\mathcal L_i$; emergency deliveries remain order-up-to.

\section*{Trip construction and route-failure probability}
Algorithm 1 (PDF page 12) defines the solution-dependent trip sets used in (8)--(9) and (16). Its printed input is a solution with values $x_{ijkt}$ for all $i,j\in\mathcal P,t\in\mathcal T$ for vehicle $k$; its output is $\mathfrak S_k$. The following reproduces its operations, including its input scope and successor expressions.
\begin{enumerate}
\item $\mathcal S\leftarrow\mathfrak S_k\leftarrow\varnothing$.
\item For $t\in\mathcal T$:
\begin{enumerate}
\item $\mathcal S_t\leftarrow\varnothing$.
\item $c\leftarrow j:\displaystyle\sum_{o'\in\mathcal O'_{kt}}\sum_{j\in\mathcal N_{kt}}x_{o'jkt}=1$.
\item While $c\notin\mathcal O''_{kt}$:
\begin{enumerate}
\item If $c\in\mathcal D$, add $\mathcal S_t$ as an element of $\mathcal S$; add $\mathcal S$ as an element of $\mathfrak S_k$; then $\mathcal S_t\leftarrow\mathcal S\leftarrow\varnothing$.
\item Else if $c\in\mathcal P$, add $c$ as an element of $\mathcal S_t$.
\item $c\leftarrow j:\displaystyle\sum_{j\in\mathcal N_{kt}}x_{cjkt}=1$.
\end{enumerate}
\item Add $\mathcal S_t$ as an element of $\mathcal S$.
\end{enumerate}
\end{enumerate}
For each trip $\mathcal S\in\mathfrak S_k$, the delivered quantity is defined by
\begin{equation}
\begin{aligned}
\Gamma_{\mathcal S}
={}&\sum_{\mathcal S_0\in\mathcal S}\sum_{s\in\mathcal S_0}(\Lambda_{s0}-I_{s0})\\
&+\sum_{t\in\mathcal T\setminus\{0\}}\sum_{\mathcal S_t\in\mathcal S}\sum_{s\in\mathcal S_t}
\left(\Lambda_{st}-\Lambda_{sm}+\sum_{h=m}^{t-1}\rho_{sh}\right),\\
&\text{where }m=\max(0,g\in\mathcal T:g<t:\exists k'\in\mathcal K:y_{sk'g}=1).
\end{aligned}\tag{8}
\end{equation}
The first summand is certain, being delivered in period zero. The second uses the most recent regular delivery, the action sequence of Definition 3, and Assumption 1. The route-failure probability is
\begin{equation}
\mathsf p^{\mathrm{RF}}_{\mathcal S,k}=\mathbb P(\Gamma_{\mathcal S}>\Omega_k).
\tag{9}
\end{equation}
Recourse is a supply-point visit. The paper describes (9) as capturing the probability of multiple failures in a trip; its printed expression is the single exceedance probability above, not an expected failure count.

\paragraph{Assumption 2: restriction for route-failure calculation only.}
The errors used to calculate route-failure probabilities are independent and identically distributed, drawn from any marginal distribution $\Phi'$. The paper stipulates
\begin{equation}
\Phi(\varepsilon)=\prod_{t\in\mathcal T}\prod_{i\in\mathcal P}\Phi'(\varepsilon_{it}),
\tag{10}
\end{equation}
where $\Phi'$ is the marginal cumulative distribution function. This iid assumption is not imposed on the general stock-out calculation. Under it, the cumulative error distribution depends on the number of summed errors, bounded above by $|\mathcal P|(|\mathcal T|-1)$. Demands realized in the last horizon period cannot be served during the horizon because delivery precedes demand realization.

\paragraph{Assumption 3.}
Supply points have unlimited inventories. No supply-inventory tracking constraint is introduced.

\section*{Collection interpretation: Appendix A}
Appendix A (PDF pages 28--29) defines $\Lambda'_{ig}$ as waste inventory after a regular collection in period $g$, distinct from the empty-space inventory used in the operative model:
\[
\Lambda_{ig}=\omega_i-\Lambda'_{ig}.
\]
The printed equivalences are
\begin{equation}
\mathbb P(\Lambda'_{ig}+\rho_{ig}\geqslant\omega_i)
=\mathbb P\bigl((\omega_i-\Lambda'_{ig})-\rho_{ig}\leqslant0\bigr)
=\mathbb P(\Lambda_{ig}-\rho_{ig}\leqslant0).
\tag{A.1}
\end{equation}
\begin{equation}
\begin{aligned}
&\mathbb P\left(\Lambda'_{ig}+\sum_{t=g}^{h}\rho_{it}\geqslant\omega_i\ \middle|\ \Lambda'_{ig}+\sum_{t=g}^{h-1}\rho_{it}<\omega_i\right)\\
&=\mathbb P\left((\omega_i-\Lambda'_{ig})-\sum_{t=g}^{h}\rho_{it}\leqslant0\ \middle|\ (\omega_i-\Lambda'_{ig})-\sum_{t=g}^{h-1}\rho_{it}>0\right)\\
&=\mathbb P\left(\Lambda_{ig}-\sum_{t=g}^{h}\rho_{it}\leqslant0\ \middle|\ \Lambda_{ig}-\sum_{t=g}^{h-1}\rho_{it}>0\right),\qquad\forall h>g.
\end{aligned}\tag{A.2}
\end{equation}
The emergency-collection cases follow as the special cases corresponding to (5)--(6). These identities justify the empty-space interpretation; they do not impose a second inventory system.

\section*{Probability evaluation definitions in Appendices B--C}
The appendices supply probability-evaluation constructions, not replacement optimization subproblems. Their simulation sample size $M$ is unrelated to the big-$M$ coefficient in (33) and (35).

\subsection*{Appendix B: stock-out probabilities}
Choose a sufficiently large $M$. For $m\in\{1,\ldots,M\}$, draw the error vector from $\Phi$:
\begin{equation}
 e^m=(e^m_{11},\ldots,e^m_{1|\mathcal T|},e^m_{21},\ldots,e^m_{|\mathcal P||\mathcal T|}).
\tag{B.1}
\end{equation}
The printed evaluation formulas are
\begin{equation}
\mathbb P(\Lambda_{ig}-\rho_{ig}\leqslant0)
=\mathbb P(\varepsilon_{ig}\geqslant\Lambda_{ig}-\mathbb E(\rho_{ig}))
=\frac{\sum_{m=1}^{M}\operatorname{IF}(e^m_{it}\geqslant\Lambda_{ig}-\mathbb E(\rho_{ig}),1,0)}{M}.
\tag{B.2}
\end{equation}
\begin{equation}
\begin{aligned}
&\mathbb P\left(\Lambda_{ig}-\sum_{t=g}^{h}\rho_{it}\leqslant0\ \middle|\ \Lambda_{ig}-\sum_{t=g}^{h-1}\rho_{it}>0\right)\\
&=\mathbb P\left(\sum_{t=g}^{h}\varepsilon_{it}\geqslant\Lambda_{ig}-\sum_{t=g}^{h}\mathbb E(\rho_{it})\ \middle|\ \sum_{t=g}^{h-1}\varepsilon_{it}<\Lambda_{ig}-\sum_{t=g}^{h-1}\mathbb E(\rho_{it})\right)\\
&=\frac{\mathbb P\left(\sum_{t=g}^{h}\varepsilon_{it}\geqslant\Lambda_{ig}-\sum_{t=g}^{h}\mathbb E(\rho_{it}),\ \sum_{t=g}^{h-1}\varepsilon_{it}<\Lambda_{ig}-\sum_{t=g}^{h-1}\mathbb E(\rho_{it})\right)}
{\mathbb P\left(\sum_{t=g}^{h-1}\varepsilon_{it}<\Lambda_{ig}-\sum_{t=g}^{h-1}\mathbb E(\rho_{it})\right)}\\
&=\frac{\displaystyle\sum_{m=1}^{M}\operatorname{IF}\left(\sum_{t=g}^{h}e^m_{it}\geqslant\Lambda_{ig}-\sum_{t=g}^{h}\mathbb E(\rho_{it})\ \mathrm{AND}\ \sum_{t=g}^{h-1}e^m_{it}<\Lambda_{ig}-\sum_{t=g}^{h-1}\mathbb E(\rho_{it}),1,0\right)}
{\displaystyle\sum_{m=1}^{M}\operatorname{IF}\left(\sum_{t=g}^{h-1}e^m_{it}<\Lambda_{ig}-\sum_{t=g}^{h-1}\mathbb E(\rho_{it}),1,0\right)},\qquad\forall h>g.
\end{aligned}\tag{B.3}
\end{equation}
$\operatorname{IF}([\text{condition}],1,0)$ equals one if the condition holds and zero otherwise. Equations (5)--(6) are evaluated as the special cases of (B.2)--(B.3) with capacity after emergency delivery. The equal signs to finite-sample frequencies and the subscript $it$ in (B.2) are retained as printed; the latter leaves $t$ unspecified in an expression about period $g$. No zero-denominator convention is supplied for (B.3).

\subsection*{Appendix C: route-failure probabilities}
\begin{equation}
\begin{aligned}
\mathbb P(\Gamma_{\mathcal S}>\Omega_k)
&=\mathbb P(\mathbb E(\Gamma_{\mathcal S})+E>\Omega_k)\\
&=\mathbb P(E>\Omega_k-\mathbb E(\Gamma_{\mathcal S})).
\end{aligned}\tag{C.1}
\end{equation}
Here plain $E$ denotes the cumulative random error, distinct from the expectation operator $\mathbb E$:
\begin{equation}
 E=\sum_{t\in\mathcal T\setminus\{0\}}\sum_{\mathcal S_t\in\mathcal S}\sum_{s\in\mathcal S_t}\sum_{h=m}^{t-1}\varepsilon_{sh},
\tag{C.2}
\end{equation}
with $m$ as in (8). Under Assumption 2, the number of iid errors is bounded by
\[
 N=|\mathcal P|(|\mathcal T|-1).
\]
For sufficiently large sample size $M$, and $m\in\{1,\ldots,M\}$, simulate
\begin{equation}
 e^m_g=\sum_{t=1}^{g}\varepsilon_{it},\qquad\forall g\in\{1,\ldots,N\},
\tag{C.3}
\end{equation}
by drawing $\varepsilon_{it}$ from $\Phi'$ for any $i\in\mathcal P$. For each $g$, derive the empirical distribution function $\Phi^{\mathrm{emp}}_g$ of $\{e^1_g,\ldots,e^M_g\}$. The appendix states
\[
\exists g\in\{1,\ldots,N\}\ \text{such that }E\sim\Phi^{\mathrm{emp}}_g,
\]
and uses these empirical distributions at runtime for (C.1). This is the printed simulation-based characterization; no exact-distribution identity or zero-summand extension is added.

\section*{Brief notes on the retained printed scope}
\begin{itemize}
\item Equation (18) equates depot departures and arrivals but does not set either sum to one, despite the routing prose and Algorithm 1 presupposing a tour origin and destination. No missing equality is inserted.
\item Equations (21) and (39) quantify over all $t\in\mathcal T$, hence include $t=u$ and refer to period $u+1$. Equation (39) also prints $o'\in\mathcal O'_{kt}$, not $\mathcal O'_{k(t+1)}$. These scopes are retained. Equation (15) likewise uses $y_{ik(u+1)}$, whereas (46) gives visit-variable domains only on $\mathcal T$. No terminal values are stipulated in the supplied formulation.
\item Definition 2 allows negative fixed $I_{i0}$, while (47) imposes $I_{it}\geqslant0$ for all $t\in\mathcal T$, including zero. Both statements are retained; no initial-inventory exception is inserted into (47).
\item The arc set excludes $i=j$, whereas several sums and (40), (46) range over all $i,j\in\mathcal N_{kt}$. The printed domains (46)--(47) also extend demand-point quantities or sets to all nodes. Equation (38) uses $q_{jkt}$ at depots, although its verbal definition is a demand-point delivery quantity. No diagonal-arc convention, depot-delivery value, or narrowed quantifier is added.
\item The prose defines stock-out as inventory becoming negative, but (3)--(6) and Appendix A use nonpositive empty-space inventory, including equality. Those printed probability boundaries are retained. Equation (2) and (B.1) index error-vector coordinates from one through $|\mathcal T|$ although $\mathcal T$ starts at zero; no reindexing rule is introduced.
\item Algorithm 1's printed input restricts $i,j$ to $\mathcal P$, although its operations use depot and supply-point arcs. Its successor notation sums over the index it seeks to select, and it does not explicitly append the remaining trip to $\mathfrak S_k$ after the outer loop. These are retained rather than repaired.
\item The statement that (55) forces a disposal visit is stronger than its algebra in zero-load cases: the equality itself does not prohibit a zero-load predecessor outside $\mathcal D$. The operative model retains (55), not an added disposal-predecessor arc constraint.
\end{itemize}
The supplied main paper and its Appendices A--C provide the dependencies reproduced here. The separately cited Markov et al. (2016a) report is not supplied and is unverified; no additional formulation or numerical assumptions from it are imported.
\end{document}